% This document was generated by an LLM.

\documentclass{essay}

\DeclareMathOperator{\sgn}{sgn}
\newcommand{\abs}[1]{\lvert #1 \rvert}

\title{The Derivative of $\abs{x}$ from First Principles}
\author{}
\date{}

\begin{document}

\maketitle

\section{Setup}

Let $f(x) = \abs{x}$. For fixed $x$ and $h \neq 0$, define the difference quotient
\[
  Q_x(h) = \frac{\abs{x+h} - \abs{x}}{h}.
\]
By definition, $f'(x) = \lim_{h \to 0} Q_x(h)$, provided the limit exists.

\begin{definition}
  Let $g$ be defined on a punctured neighborhood of $0$. Then
  $\lim_{h \to 0} g(h) = L$ means: for every $\varepsilon > 0$ there exists
  $\delta > 0$ such that
  \[
    0 < \abs{h} < \delta \implies \abs{g(h) - L} < \varepsilon.
  \]
\end{definition}

Two features of this definition are used throughout.

\begin{enumerate}
  \item The number $\delta$ is chosen by the prover. If some $\delta$ works,
        every smaller positive $\delta$ also works, since it only restricts
        $h$ further.
  \item The condition says nothing about $h$ with $\abs{h} \geq \delta$.
        Values of $g$ outside the chosen neighborhood never enter the
        definition. A limit is a local property.
\end{enumerate}

\begin{lemma}[Local agreement]\label{lem:local}
  Suppose $g(h) = k(h)$ for all $h$ with $0 < \abs{h} < r$, for some $r > 0$.
  Then $\lim_{h \to 0} g(h) = L$ if and only if $\lim_{h \to 0} k(h) = L$.
\end{lemma}

\begin{proof}
  Suppose $\lim_{h\to 0} k(h) = L$. Given $\varepsilon > 0$, let $\delta$ be
  the number provided for $k$, and set $\delta' = \min(\delta, r)$. If
  $0 < \abs{h} < \delta'$, then $g(h) = k(h)$ and $\abs{k(h) - L} < \varepsilon$,
  so $\abs{g(h) - L} < \varepsilon$. The converse is symmetric.
\end{proof}

Every derivation below reduces the problem to this lemma: one shows that
$Q_x(h)$ agrees with a simpler expression for all sufficiently small $h \neq 0$,
and then takes the limit of the simpler expression.

\section{Method 1: Case analysis}

The absolute value is defined piecewise:
$\abs{a} = a$ for $a \geq 0$ and $\abs{a} = -a$ for $a \leq 0$.
The task is to show that, for $h$ small enough, $x + h$ has the same sign as $x$,
so that a single piece of the definition applies to both $\abs{x+h}$ and $\abs{x}$.

\subsection{The case $x > 0$}

\begin{proposition}
  If $x > 0$, then $f'(x) = 1$.
\end{proposition}

\begin{proof}
  Take $0 < \abs{h} < x$. Since $h \geq -\abs{h}$, with equality exactly when
  $h < 0$,
  \[
    x + h \geq x - \abs{h} > 0.
  \]
  The first inequality cannot be made strict: for $h < 0$ it is an equality.
  Strictness of the conclusion comes from the second inequality, which follows
  from $\abs{h} < x$; a chain $a \geq b > c$ gives $a > c$.

  Hence $\abs{x+h} = x + h$ and $\abs{x} = x$, so
  \[
    Q_x(h) = \frac{(x+h) - x}{h} = \frac{h}{h} = 1.
  \]
  Given any $\varepsilon > 0$, choose $\delta = x$. Then
  $0 < \abs{h} < \delta$ implies $\abs{Q_x(h) - 1} = 0 < \varepsilon$.
  Thus $f'(x) = 1$.
\end{proof}

Here a single $\delta$ works for every $\varepsilon$. This happens only because
$Q_x$ is exactly constant near $0$.

\subsection{The case $x < 0$}

\begin{proposition}
  If $x < 0$, then $f'(x) = -1$.
\end{proposition}

\begin{proof}
  Take $0 < \abs{h} < \abs{x} = -x$. Since $h \leq \abs{h}$, with equality
  exactly when $h > 0$,
  \[
    x + h \leq x + \abs{h} < 0.
  \]
  Hence $\abs{x+h} = -(x+h)$ and $\abs{x} = -x$, so
  \[
    Q_x(h) = \frac{-(x+h) - (-x)}{h} = \frac{-h}{h} = -1.
  \]
  Choosing $\delta = -x$ completes the proof as before.
\end{proof}

\subsection{The case $x = 0$}

For $h \neq 0$,
\[
  Q_0(h) = \frac{\abs{h} - 0}{h} = \frac{\abs{h}}{h} =
  \begin{cases}
    1,  & h > 0, \\
    -1, & h < 0.
  \end{cases}
\]
The right-hand limit is $1$ and the left-hand limit is $-1$.

\begin{proposition}
  $f$ is not differentiable at $0$.
\end{proposition}

\begin{proof}
  Suppose $\lim_{h \to 0} Q_0(h) = L$ for some real $L$. Take $\varepsilon = 1$
  and let $\delta > 0$ be the corresponding number. Both $h_1 = \delta/2$ and
  $h_2 = -\delta/2$ satisfy $0 < \abs{h} < \delta$. From $h_1$:
  $\abs{1 - L} < 1$, so $0 < L < 2$. From $h_2$: $\abs{-1 - L} < 1$, so
  $-2 < L < 0$. These are incompatible.
\end{proof}

\subsection{Remarks on the choice of neighborhood}

\begin{remark}[Why not arbitrary $h$]
  The identity $Q_x(h) = 1$ for $x > 0$ is false for large negative $h$.
  For example, with $x = 1$ and $h = -3$,
  \[
    Q_1(-3) = \frac{\abs{-2} - \abs{1}}{-3} = \frac{2 - 1}{-3} = -\frac{1}{3}.
  \]
  Once $h < -x$, the sign of $x + h$ flips and the formula
  $\abs{x+h} = x + h$ no longer applies. This does not affect the limit:
  by feature 2 of the definition and Lemma~\ref{lem:local}, only the values of
  $Q_x$ on some punctured neighborhood of $0$ matter, and the restriction
  $0 < \abs{h} < x$ is precisely the choice $\delta = x$ that the prover is
  entitled to make.
\end{remark}

\begin{remark}[Non-strict neighborhood]
  One may instead take $0 < \abs{h} \leq x$. Then $h = -x$ is allowed and gives
  $x + h = 0$, so the claim $x + h > 0$ fails. What survives is
  \[
    x + h \geq x - \abs{h} \geq 0,
  \]
  which suffices, because $\abs{a} = a$ holds for all $a \geq 0$. Directly,
  \[
    Q_x(-x) = \frac{0 - x}{-x} = 1.
  \]
  Similarly for $x < 0$ with $0 < \abs{h} \leq -x$: at $h = -x$ we have
  $x + h = 0$, and $\abs{a} = -a$ holds for all $a \leq 0$, so $Q_x(-x) = -1$.

  Nothing is gained. Any smaller neighborhood, such as
  $0 < \abs{h} < x/2$, works equally well. The strict version is conventional
  because it matches the form $0 < \abs{h} < \delta$ in the definition and
  avoids the boundary case $x + h = 0$.
\end{remark}

\begin{remark}[The largest set]
  For $x > 0$, the set of all $h \neq 0$ with $Q_x(h) = 1$ is not symmetric:
  it is $\{h \neq 0 : h \geq -x\}$, since positive $h$ of any size keeps
  $x + h > 0$. The symmetric neighborhood is chosen only because it fits the
  definition of the limit directly.
\end{remark}

\section{Method 2: Rationalization}

\subsection{Motivation}

The identity $\abs{a} = \sqrt{a^2}$ rewrites the difference quotient as
\[
  Q_x(h) = \frac{\sqrt{(x+h)^2} - \sqrt{x^2}}{h},
\]
which has the structure of the first-principles problem for $\sqrt{x}$: a
difference of two roots over $h$. The standard technique there is to multiply
numerator and denominator by the sum of the roots, so that the identity
$(A - B)(A + B) = A^2 - B^2$ removes the roots. The same technique applies with
$\abs{\cdot}$ in place of $\sqrt{\cdot}$.

The quotient as written has the form $0/0$ in the limit and cannot be evaluated
by substituting $h = 0$. The goal is to rewrite it until a factor of $h$ cancels.

\subsection{Derivation}

Fix $x \neq 0$ and $h \neq 0$. Write $A = \abs{x+h}$ and $B = \abs{x}$.

\paragraph{Step 1: Multiply by $1$ in a useful form.}
\[
  \frac{A - B}{h} = \frac{A - B}{h} \cdot \frac{A + B}{A + B}.
\]
This requires $A + B \neq 0$. Since $B = \abs{x} > 0$ and $A \geq 0$, we have
$A + B \geq \abs{x} > 0$. This is the first use of $x \neq 0$.

\paragraph{Step 2: Difference of squares.}
Since $(A - B)(A + B) = A^2 - B^2$,
\[
  \frac{A - B}{h} = \frac{A^2 - B^2}{h\,(A + B)}.
\]

\paragraph{Step 3: Remove the absolute values in the numerator.}
For every real $a$, $\abs{a}^2 = a^2$. Hence
\[
  A^2 = (x+h)^2, \qquad B^2 = x^2.
\]
This is the essential step. Before it, the numerator contains absolute values,
which force a case split. After it, the numerator is a polynomial. The absolute
values remain in the denominator, where they cause no difficulty.

\paragraph{Step 4: Expand.}
\[
  (x+h)^2 - x^2 = (x^2 + 2xh + h^2) - x^2 = 2xh + h^2 = h(2x + h).
\]

\paragraph{Step 5: Cancel $h$.}
\[
  Q_x(h) = \frac{h(2x + h)}{h\,(\abs{x+h} + \abs{x})}
         = \frac{2x + h}{\abs{x+h} + \abs{x}}.
\]
The cancellation is valid because $h \neq 0$. The two expressions agree for all
$h \neq 0$, which is a punctured neighborhood of $0$, so by
Lemma~\ref{lem:local} they have the same limit.

\paragraph{Step 6: Take the limit.}
The numerator satisfies $2x + h \to 2x$. For the denominator, the reverse
triangle inequality gives
\[
  \bigl\lvert\, \abs{x+h} - \abs{x} \,\bigr\rvert \leq \abs{h} \to 0,
\]
so $\abs{x+h} \to \abs{x}$ and $\abs{x+h} + \abs{x} \to 2\abs{x}$. The quotient
rule for limits requires the limit of the denominator to be nonzero; it is
$2\abs{x} > 0$, which is the second use of $x \neq 0$. Therefore
\[
  f'(x) = \lim_{h \to 0} \frac{2x + h}{\abs{x+h} + \abs{x}}
        = \frac{2x}{2\abs{x}} = \frac{x}{\abs{x}}.
\]

\paragraph{Step 7: Read off the sign.}
$x / \abs{x}$ equals $1$ for $x > 0$ and $-1$ for $x < 0$, in agreement with
Method 1.

\subsection{Where each hypothesis is used}

\begin{center}
  \begin{tabular}{lll}
    Step & Requirement & Reason \\
    \hline
    1 & $\abs{x+h} + \abs{x} \neq 0$ & multiplying by $(A+B)/(A+B)$ \\
    3 & none & $\abs{a}^2 = a^2$ for all real $a$ \\
    5 & $h \neq 0$ & cancelling $h$; always true inside a limit \\
    6 & $\abs{x} \neq 0$ & quotient rule needs a nonzero denominator limit \\
  \end{tabular}
\end{center}

\subsection{What happens at $x = 0$}

With $x = 0$, Step 1 requires $\abs{h} \neq 0$, which holds. Steps 2--5 give
\[
  Q_0(h) = \frac{2 \cdot 0 + h}{\abs{h} + 0} = \frac{h}{\abs{h}}.
\]
Step 6 fails: the denominator tends to $0$, the quotient rule does not apply,
and $h / \abs{h} = \pm 1$ according to the sign of $h$. This is the original
quotient $\abs{h}/h$ again, which has no limit. The algebra locates the failure
exactly.

\section{Method 3: The chain rule}

This method is not from first principles, since it presupposes the derivative
of $\sqrt{u}$, but it shows the same structure from another angle.

Write $f(x) = \sqrt{x^2}$, the composition of $x \mapsto x^2$ with
$u \mapsto \sqrt{u}$. For $x \neq 0$, the outer function is differentiable at
$u = x^2 > 0$ with derivative $1/(2\sqrt{u})$. The chain rule gives
\[
  f'(x) = \frac{1}{2\sqrt{x^2}} \cdot 2x = \frac{x}{\abs{x}}.
\]
The derivative of $\sqrt{u}$ for $u > 0$ is itself obtained by rationalization,
so this method is Method 2 packaged through a theorem.

At $x = 0$ the outer function is evaluated at $u = 0$, where $\sqrt{u}$ is not
differentiable, so the chain rule does not apply. This proves nothing about
$f'(0)$: the chain rule is a sufficient condition, and the failure of a
sufficient condition does not establish non-differentiability. For example,
$x^2 = \sqrt{x^4}$ has the same form, with $\sqrt{u}$ again evaluated at
$u = 0$ when $x = 0$, yet $x^2$ is differentiable at $0$. Non-differentiability
of $\abs{x}$ at $0$ must be shown directly, as in Method 1.

\section{Method 4: Locality of the derivative}

Method 1 can be stated as a general principle.

\begin{proposition}
  Suppose $f(t) = g(t)$ for all $t$ in an open interval $(x - r, x + r)$, and
  $g$ is differentiable at $x$. Then $f$ is differentiable at $x$ and
  $f'(x) = g'(x)$.
\end{proposition}

\begin{proof}
  For $0 < \abs{h} < r$, both $x$ and $x + h$ lie in the interval, so
  \[
    \frac{f(x+h) - f(x)}{h} = \frac{g(x+h) - g(x)}{h}.
  \]
  The difference quotients agree on a punctured neighborhood of $0$, and
  Lemma~\ref{lem:local} gives the result.
\end{proof}

For $x > 0$, $\abs{t} = t$ on $(0, 2x) = (x - x, x + x)$, so
$f'(x) = \frac{d}{dt} t = 1$. For $x < 0$, $\abs{t} = -t$ on $(2x, 0)$, so
$f'(x) = -1$. The radius $r = \abs{x}$ is the same restriction
$0 < \abs{h} < \abs{x}$ used in Method 1.

At $x = 0$ no open interval around $0$ lies within a single piece. As with the
chain rule, this observation alone does not prove non-differentiability. The
one-sided version of the proposition does apply: $\abs{t} = t$ on $[0, \infty)$
and $\abs{t} = -t$ on $(-\infty, 0]$, so the right derivative at $0$ is $1$ and
the left derivative is $-1$. Since they differ, the two-sided derivative does
not exist.

\section{Result}

\[
  \frac{d}{dx}\abs{x} = \frac{x}{\abs{x}} = \sgn(x), \qquad x \neq 0,
\]
and the derivative does not exist at $x = 0$. The function $\abs{x}$ is
continuous at $0$ but not differentiable there.

\begin{remark}
  The one-sided derivatives at $0$ exist:
  \[
    f'_+(0) = \lim_{h \to 0^+} \frac{\abs{h}}{h} = 1, \qquad
    f'_-(0) = \lim_{h \to 0^-} \frac{\abs{h}}{h} = -1.
  \]
\end{remark}

\begin{remark}
  In convex analysis the subdifferential of $\abs{x}$ at $0$ is the interval
  $[-1, 1]$, the set of slopes of lines through the origin lying below the
  graph. This is a different object from the derivative.
\end{remark}

\begin{remark}
  The formula $\frac{d}{dx}\abs{x} = \sgn(x)$ holds only on $x \neq 0$. It is an
  error to extend it to $0$ by the convention $\sgn(0) = 0$: the derivative at
  $0$ does not exist, so it has no value, $0$ or otherwise.
\end{remark}

\section{Self-check questions}

\paragraph{Q1.} Why may the derivation for $x > 0$ restrict to
$0 < \abs{h} < x$?

\noindent\textit{Answer.} The definition of the limit lets the prover choose
$\delta$, and only values with $0 < \abs{h} < \delta$ are examined. Choosing
$\delta = x$ is legitimate; values of $Q_x$ outside this neighborhood are
irrelevant.

\paragraph{Q2.} Compute $Q_1(-3)$. What does the result show?

\noindent\textit{Answer.} $Q_1(-3) = -1/3$. The identity $Q_x(h) = 1$ fails
for large negative $h$, which is why the neighborhood must be restricted.

\paragraph{Q3.} Does the derivation for $x > 0$ still work with
$0 < \abs{h} \leq x$?

\noindent\textit{Answer.} Yes. At $h = -x$, $x + h = 0$ and $\abs{0} = 0 = x + h$,
so $Q_x(-x) = 1$. The justification uses $x + h \geq 0$ instead of
$x + h > 0$.

\paragraph{Q4.} In the rationalization method, which step removes the absolute
values, and why does it work?

\noindent\textit{Answer.} Step 3, using $\abs{a}^2 = a^2$ for all real $a$.
Squaring discards the sign, turning the numerator into a polynomial.

\paragraph{Q5.} Where is $x \neq 0$ used in the rationalization method?

\noindent\textit{Answer.} In Step 1, to guarantee $\abs{x+h} + \abs{x} \neq 0$,
and in Step 6, so that the denominator tends to $2\abs{x} \neq 0$ and the
quotient rule for limits applies.

\paragraph{Q6.} Why is cancelling $h$ in Step 5 valid?

\noindent\textit{Answer.} In the limit $h \to 0$, $h$ is never $0$. The two
expressions agree on the punctured neighborhood $h \neq 0$, so by local
agreement they have the same limit.

\paragraph{Q7.} The chain rule does not apply to $\sqrt{x^2}$ at $x = 0$. Does
this prove $\abs{x}$ is not differentiable at $0$?

\noindent\textit{Answer.} No. The chain rule is a sufficient condition. The
function $\sqrt{x^4} = x^2$ has the same obstruction and is differentiable at
$0$. Non-differentiability requires the direct one-sided argument.

\paragraph{Q8.} Is $\frac{d}{dx}\abs{x}$ equal to $0$ at $x = 0$?

\noindent\textit{Answer.} No. The derivative does not exist at $0$. The
one-sided derivatives are $1$ and $-1$.

\end{document}
